\chapter{Comparison of Optimized Designs with Laminar Baselines}
\label{app:laminar_turbulent_topology_comparison}

\providecommand{\validationplaceholder}[1]{%
	\fbox{\parbox[c][0.22\textheight][c]{0.88\linewidth}{\centering \textbf{Placeholder}\\[0.5em]#1}}%
}

While Chapter~\ref{ch:benchmark_results} successfully verifies the computational accuracy of the fully coupled FEniCS-SA framework against commercial CFD solvers, an essential engineering justification underpins this methodology: does optimizing a fluid manifold under turbulent conditions yield a physically superior design compared to a standard laminar baseline? Because continuous turbulent adjoint pipelines demand substantially more computational resources, tighter iterative coupling, and rigorous stabilization compared to their laminar equivalents, this mathematical complexity must translate into a measurable reduction in hydrodynamic losses. This appendix quantifies that performance gap.

\section{Methodology and Evaluation Pipeline}
\label{sec:app_lam_vs_turb_methodology}
To conduct a rigorous comparative analysis, a representative routing domain (the [Insert Benchmark] case) was optimized twice utilizing the custom FEniCS framework. The first design iteration was generated using the steady-state, incompressible laminar Navier-Stokes equations, effectively blinding the optimizer to turbulent diffusion and Reynolds stresses. The second iteration utilized the fully coupled Spalart-Allmaras methodology developed in Chapter~\ref{ch:topology_optimization_methodology}. 

Both resulting continuous design fields were subjected to the identical geometry extraction and thresholding procedure. To evaluate their true operational performance, the sharp CAD manifolds were imported into Ansys Fluent, meshed with wall-resolved inflation layers ($y^+ \le 1$), and subjected to identical turbulent boundary conditions utilizing the SST $k$-$\omega$ turbulence model. The inlets were prescribed a uniform velocity profile of [Insert Velocity] m/s and a turbulence intensity of [Insert TI]\%, corresponding to a highly convective operational Reynolds number of $Re = \text{[Insert Re]}$.

\section{Morphological and Hydrodynamic Comparison}
\label{sec:app_lam_vs_turb_morphology}
The geometric discrepancies between the laminar and turbulent topologies are stark, directly reflecting the distinct physical phenomena captured by the respective adjoint sensitivities. Figure~\ref{fig:app_velocity_comparison} illustrates the velocity magnitude contours for both geometries evaluated under the high-Reynolds boundary conditions.

\begin{figure}[!ht]
	\centering
	\begin{subfigure}[b]{0.47\textwidth}
		\centering
		\validationplaceholder{Extracted Laminar Topology and Velocity Field (Ansys SST)}
		\caption{Laminar-optimized design evaluated under turbulent flow.}
		\label{fig:app_laminar_vel}
	\end{subfigure}
	\hfill
	\begin{subfigure}[b]{0.47\textwidth}
		\centering
		\validationplaceholder{Extracted Turbulent Topology and Velocity Field (Ansys SST)}
		\caption{Turbulent-optimized design (SA) evaluated under turbulent flow.}
		\label{fig:app_turbulent_vel}
	\end{subfigure}
	\caption{Velocity magnitude contours of the laminar and turbulent designs, evaluated under identical high-Reynolds boundary conditions in Ansys Fluent (SST $k$-$\omega$).}
	\label{fig:app_velocity_comparison}
\end{figure}

As observed in Figure~\ref{fig:app_velocity_comparison}(a), the laminar optimizer favors abrupt expansions and sharp internal routing. Because the laminar adjoint formulation lacks the convective momentum transport and severe boundary layer mechanics associated with inertial flows, it distributes structural material strictly to minimize bulk viscous friction. Consequently, when this laminar geometry operates in a turbulent regime, the severe adverse pressure gradients trigger massive boundary layer separation, generating highly dissipative recirculation zones. 

Conversely, Figure~\ref{fig:app_velocity_comparison}(b) demonstrates the mathematical response of the turbulent optimizer. The SA-driven framework actively suppresses separation by generating heavily streamlined, gradual geometric transitions. It trades a minor increase in wetted surface area (and subsequent skin friction) for a drastic reduction in form drag and turbulent wake generation.

\section{Quantitative Performance Analysis}
\label{sec:app_lam_vs_turb_metrics}
The physical differences in flow morphology translate directly into macroscopic thermodynamic performance. Table~\ref{tab:app_laminar_turbulent_metrics} summarizes the power dissipation and peak turbulence metrics evaluated by the Ansys SST $k$-$\omega$ solver.

\begin{table}[!ht]
	\centering
	\footnotesize
	\caption{Quantitative performance metrics of the laminar and turbulent optimized topologies evaluated under identical turbulent boundary conditions (Ansys SST $k$-$\omega$).}
	\label{tab:app_laminar_turbulent_metrics}
	\renewcommand{\arraystretch}{1.2}
	\begin{tabularx}{\textwidth}{l c c X}
		\hline
		\textbf{Design Origin} & \textbf{Total Pressure Drop ($\Delta p$) [Pa]} & \textbf{Peak Eddy Viscosity Ratio} & \textbf{Relative Performance} \\
		\hline
		Laminar TO Formulation & [TBD] & [TBD] & Baseline \\
		Turbulent TO Formulation (SA) & [TBD] & [TBD] & \textbf{-[TBD]\%} \\
		\hline
	\end{tabularx}
\end{table}

The turbulent-optimized manifold exhibits a significant [Insert \%] reduction in total pressure drop compared to the laminar baseline. This severe performance penalty in the laminar geometry is heavily dominated by the artificial turbulent kinetic energy generated within its flow separation zones. Ultimately, this independent CFD re-analysis confirms that extending topology optimization into the turbulent regime is not merely an academic exercise, but a strict engineering requirement to ensure the manufacturability and efficiency of high-Reynolds-number fluid systems.